\documentclass[11pt]{article}
\usepackage[T1]{fontenc}
\usepackage{textcomp}
\usepackage[margin=0.75in]{geometry}
\usepackage{amsmath,amssymb,amsthm}
\usepackage{array,booktabs}
\usepackage{hyperref}
\setlength{\parindent}{0pt}
\newtheorem{theorem}{Theorem}
\theoremstyle{definition}
\newtheorem{definition}{Definition}
\title{Flow Proof Artifact: RingHom}
\author{Generated by \texttt{flow doc proof}}
\date{Flow Proof Book}
\begin{document}
\maketitle
Ring homomorphism axioms.\\[0.5em]
\textbf{Source.} dummit-foote — *Abstract Algebra*, §7.3\\[0.5em]
\bigskip
\noindent{\large \textbf{Definition 1.} \textit{A ring homomorphism preserves addition} \hfill RingHom \cdot addition \cdot sums map to sums}\par
\smallskip\noindent\textit{Coordinate: RingHom \cdot addition \cdot sums map to sums}\par
\smallskip\noindent\textit{Source: dummit-foote}\par
\medskip
\noindent\textbf{Goal.} A ring homomorphism preserves addition\par
\noindent$\displaystyle \forall f \in RingHom \forall a \in Ring \forall b \in Ring\quad f(a + b) = f(a) + f(b)$\par
\medskip
\noindent
\begin{tabular}{@{} >{\raggedright\arraybackslash}p{0.06\textwidth} >{\raggedright\arraybackslash}p{0.47\textwidth} >{\raggedleft\arraybackslash}p{0.41\textwidth} @{}}
\textbf{\#} & \textbf{Proof} & \textbf{Mathematics} \\
\midrule
\textbf{1.} & We stipulate sums map to sums for addition on RingHom — this is a definition, not a derived fact. & \\[0.45em]
\textbf{2.} & This follows directly from the definition: f(a plus b) equals f(a) plus f(b). Hence proven. & $\displaystyle f(a + b) = f(a) + f(b)$ \\[0.45em]
\bottomrule
\end{tabular}
\label{thm:RingHom-addition-sums_map_to_sums}
\par\medskip\hrule
\bigskip
\noindent{\large \textbf{Definition 2.} \textit{A ring homomorphism preserves multiplication} \hfill RingHom \cdot multiplication \cdot products map to products}\par
\smallskip\noindent\textit{Coordinate: RingHom \cdot multiplication \cdot products map to products}\par
\smallskip\noindent\textit{Source: dummit-foote}\par
\medskip
\noindent\textbf{Goal.} A ring homomorphism preserves multiplication\par
\noindent$\displaystyle \forall f \in RingHom \forall a \in Ring \forall b \in Ring\quad f(a \cdot b) = f(a) * f(b)$\par
\medskip
\noindent
\begin{tabular}{@{} >{\raggedright\arraybackslash}p{0.06\textwidth} >{\raggedright\arraybackslash}p{0.47\textwidth} >{\raggedleft\arraybackslash}p{0.41\textwidth} @{}}
\textbf{\#} & \textbf{Proof} & \textbf{Mathematics} \\
\midrule
\textbf{1.} & We stipulate products map to products for multiplication on RingHom — this is a definition, not a derived fact. & \\[0.45em]
\textbf{2.} & This follows directly from the definition: f(a times b) equals f(a) times f(b). Hence proven. & $\displaystyle f(a \cdot b) = f(a) * f(b)$ \\[0.45em]
\bottomrule
\end{tabular}
\label{thm:RingHom-multiplication-products_map_to_products}
\par\medskip\hrule
\end{document}
